% Owned by the skeleton thread.
\section{Introduction}
\label{sec:intro}

Full waveform inversion (FWI) estimates the Earth's elastic properties by fitting whole
recorded seismograms with simulated ones \citep{tarantola1984,mora1987,fichtner2011}. Each
iteration simulates the wavefield of every source in the current model, sends the misfit
between simulated and recorded traces back through the model as an adjoint wavefield, and
correlates the two to get the misfit's gradient with respect to the model
\citep{tromp2005}. Everything an inversion concludes therefore rests on the simulator: an
error in the simulated wavefield is indistinguishable, to the inversion, from an error in
the model, and a gradient that is not the gradient of what the simulator computes sends the
optimiser the wrong way.

\vr{} is a simulator written for that role. It solves the isotropic elastic wave equation
in three dimensions (and, on a grid one node across, in two) in its velocity--stress form on
a staggered grid \citep{virieux1986,levander1988,graves1996}, with convolutional perfectly
matched layers on the sides \citep{komatitsch2007} and a free surface on top. Around that
forward solver it has what an inversion needs: the adjoint run, gradients with respect to
$\vp$, $\vs$ and $\rho$, several misfit functions, Born and Hessian-vector products, and a
preconditioned L-BFGS inversion. It is written in Rust, runs on CPUs (and its time step on an
NVIDIA GPU), in a web browser as WebAssembly, and on a remote machine behind a small server.

This report is the case that the solver is correct, and a manual for using it. It sets out
to do four things.

\begin{description}
  \item[State the method.] \Cref{sec:theory} derives the equations the solver integrates and
    the adjoint equations and gradients built on them; \cref{sec:numerics} gives the discrete
    scheme, from the stencils and their stability limit to the absorbing layers, the free
    surface, and how sources, receivers and material interfaces are put on the grid.
  \item[Show that it is right.] \Cref{sec:verification} compares whole simulations with
    closed-form solutions of the same problems, and checks the properties an inversion
    depends on: reciprocity, the adjoint as the exact transpose of the forward run, and
    gradients against finite differences and exact sensitivity kernels.
  \item[Place it beside published work.] \Cref{sec:benchmarks} runs problems from the
    literature and compares the results, and the cost, with what others report.
  \item[Teach its use.] \Cref{sec:usage} is a tutorial that starts from an empty terminal.
\end{description}

\paragraph{Every number can be traced.} The solver is still being developed, so a report of
its accuracy written by hand would be out of date at the next commit. This one is generated.
Each table, figure and quoted number comes from an \emph{experiment}: a list of commands,
kept in the repository, whose reduced results are committed together with the commit they
were run at, the machine, and the time they took. Captions end with the experiment's name
and commit; \cref{sec:reproducibility} lists the commands of every experiment and says which,
if any, were made from code that has changed since. Rebuilding the report after a change to
the solver shows, as a difference of numbers, exactly what the change did to its accuracy.

\pending{a summary table of the headline results, drawn from the verification and benchmark
  experiments' values once those exist; what the report does not cover (anisotropy,
  attenuation, topography)}
